\chapter{Complément : la géométrie des probabilités}
\section*{Exercice 28.1 -- Moyennes de groupes}
Sur quatre issues équiprobables, $X$ vaut $1,3,4,8$. On connaît seulement si l'issue est dans la première paire ou la seconde. Calculer son espérance conditionnelle.

\section*{Exercice 28.2 -- Variance totale}
Pour les données précédentes, vérifier la décomposition de la variance.

\section*{Exercice 28.3 -- Erreur d'un prédicteur}
Avec les mêmes groupes, comparer $Y=(2,2,6,6)$ à $Z=(1,1,5,5)$ pour prédire $X=(1,3,4,8)$.

\section*{Exercice 28.4 -- Corrélation extrême}
Si $Y=-3X+2$ avec $\operatorname{Var}(X)>0$, calculer la corrélation.

\section*{Exercice 28.5 -- Jensen}
Prouver $\mathbb E(e^X)\ge e^{\mathbb E(X)}$ pour une variable finie.

\section*{Exercice 28.6 -- Information complète}
Que devient l'espérance conditionnelle si la partition sépare toutes les issues de probabilité positive ? Et si elle ne contient que $\Omega$ ?
